\section{References}
\label{sec:aquinas-references}

\subsection{Life sources, documents, and text controls}

\begin{itemize}
\item Alarcón, Enrique, ed.  \emph{Fontes vitae Thomae Aquinatis},
\sourceurl{https://www.corpusthomisticum.org/ibfonvit.html}{Corpus Thomisticum
source map}; accessed 2026-07-20.  Used to identify contemporary documents,
early notices, Tocco, the canonization inquiries, Gui, Calo, and their editions.
The electronic presentation is rights-reserved; this publication links and
paraphrases rather than reproducing it.

\item Prümmer, D. M., and M.-H. Laurent, eds.  \emph{Fontes vitae S. Thomae
Aquinatis notis historicis et criticis illustrati}, six fascicles (Toulouse and
Saint-Maximin, 1911--1937),
\sourceurl{https://archive.org/details/fontesvitaesthom00pr}{combined scan};
accessed 2026-07-20.  Includes Peter Calo, William of Tocco, Bernard Gui, the
Naples 1319 and Fossanova 1321 inquiries, canonization narratives, and
documentary materials.  Page images, not OCR alone, control located claims;
the collection's editorial-rights history is not assumed uniform.

\item William of Tocco.  \emph{Ystoria sancti Thome de Aquino}, composed and
revised during promotion of the canonization cause, especially ch.~18 on works
and the Corpus Christi commission,
\sourceurl{https://www.corpusthomisticum.org/lgtocc18.html}{Latin working text};
and Claire Le Brun-Gouanvic, ed., \emph{Ystoria sancti Thome de Aquino de
Guillaume de Tocco (1323)} (Toronto: PIMS, 1996),
\sourceurl{https://pims.ca/publication/isbn-978-0-88844-127-0/}{edition record};
accessed 2026-07-20.  The modern edition remains protected.

\item \emph{Liber de inquisitione super vita et conversatione et miraculis
fratris Thomae de Aquino}, Naples, 1319, and \emph{Inquisitio de miraculis},
Fossanova, 1321, in \emph{Fontes vitae}, fascicles 4--5.  Used speaker by
speaker for companion testimony, transmitted reports, cult, and miracles; an
official canonization record is not treated as one eyewitness voice.

\item Ptolemy of Lucca, \emph{Historia ecclesiastica nova} 22.20--25, 39;
23.8--16; Bernard Gui, \emph{Vita Sancti Thomae Aquinatis}; and Peter Calo,
\emph{Vita S. Thomae Aquinatis}, editions mapped in the
\sourceurl{https://www.corpusthomisticum.org/ibfonvit.html}{\emph{Fontes vitae}
catalog}.  Used as distinct early witnesses only where their dependence and
genre permit.

\item Alexander IV's letter to the chancellor of Paris concerning Thomas's
licence to teach (fascicle 6, document XI, pp.~544--545); the 1259 Valenciennes
study ordinances (document XV, pp.~559--561); Charles I's 15 October 1272
allowance for Thomas's Naples lectures (document XXVIII, pp.~579--580); and
other family, university, royal, and papal documents in \emph{Fontes vitae},
fascicle 6.  These control institutions and dates, not the motives assigned by
later lives.

\item Urban IV.  \emph{Transiturus de mundo}, 11 August 1264,
\sourceurl{https://www.vatican.va/content/urbanus-iv/la/documents/bulla-transiturus-de-mundo-11-aug-1264.html}{Holy See Latin}; accessed
2026-07-20.  The bull controls extension of the feast and its enclosed proper
texts, not the identity of their author or the transmission of every received
office, Mass, hymn, and prayer.

\item Denifle, Heinrich, and Émile Châtelain, eds.  \emph{Chartularium
Universitatis Parisiensis}, vols.~1--2 (Paris: Delalain, 1889--1891): vol.~1,
no.~432, pp.~486--487 (1270); no.~473, pp.~543--558 (Paris, 1277); pp.~558--560
(Oxford, 1277); and vol.~2, no.~838, pp.~280--282 (1325),
\sourceurl{https://archive.org/details/chartulariumuniv01univuoft}{volume 1} and
\sourceurl{https://archive.org/details/chartulariumuniv02univuoft}{volume 2};
accessed 2026-07-20.  The acts control their own authorities, articles, dates,
and stated effects; they do not name Thomas in the Paris lists or define one
universal condemnation of Thomism.

\item Piché, David, with Claude Lafleur.  \emph{La condamnation parisienne de
1277: Nouvelle édition du texte latin, traduction, introduction et
commentaire} (Paris: Vrin, 1999),
\sourceurl{https://www.vrin.fr/livre/9782711614165/la-condamnation-parisienne-de-1277}{publisher record}; accessed 2026-07-20.  Used as the modern critical control
for the 1277 Paris act; protected edition, translation, and commentary are
paraphrased rather than reproduced.

\item Silva, José Filipe.  ``Robert Kilwardby,'' \emph{Stanford Encyclopedia
of Philosophy}, substantive revision 30 December 2025,
\sourceurl{https://plato.stanford.edu/entries/robert-kilwardby/}{online article};
accessed 2026-07-20.  Used to distinguish the Oxford prohibitions from the
Paris condemnation and to bound their relation to Thomistic positions.

\item Thomas Aquinas.  \emph{Opera omnia iussu Leonis XIII P. M. edita},
continuing Leonine edition; volume and authenticity orientation through the
\sourceurl{https://www.corpusthomisticum.org/reoptedi.html}{Corpus Thomisticum
best-editions index} and
\sourceurl{https://www.corpusthomisticum.org/reopauth.html}{authentication
catalog}; accessed 2026-07-20.  Critical editions and apparatus retain their
respective rights.

\item Thomas Aquinas.  \emph{Opera omnia}, searchable Latin working corpus,
\sourceurl{https://www.corpusthomisticum.org/iopera.html}{Corpus Thomisticum};
accessed 2026-07-20.  Used at exact work loci and with the site's current
probable/doubtful classifications for eucharistic texts.  Medieval wording is
public domain; the electronic edition and presentation are not relicensed.

\item Thomas Aquinas.  \emph{Summa theologiae}, especially I, qq.~1--2,
75--76, 92--93; I--II, qq.~1--5, 55--67, 90--108; II--II, qq.~10--11,
57, 101--104, 166, 179--189; and III, qq.~1--90.  Latin Leonine-oriented text
through Corpus Thomisticum; historical English Dominican translation through
\sourceurl{https://www.newadvent.org/summa/}{New Advent}; accessed 2026-07-20.
The English is a public-domain research aid, not the controlling text.

\item Thomas Aquinas.  \emph{Scriptum super Sententiis}; \emph{De ente et
essentia}; \emph{Summa contra gentiles}; \emph{Contra impugnantes Dei cultum
et religionem}; \emph{De perfectione spiritualis vitae}; \emph{Contra
doctrinam retrahentium a religione}; \emph{Catena aurea}; \emph{Contra errores
Graecorum}; \emph{De aeternitate mundi}; \emph{De unitate intellectus contra
Averroistas}; disputed questions, quodlibets, biblical and Aristotelian
commentaries.  Exact authentication and chronology follow the catalogs above;
no single translation was reproduced.

\item Thomas Aquinas.  \emph{Epistola ad ducissam Brabantiae}, traditionally
\emph{De regimine Iudaeorum},
\sourceurl{https://www.corpusthomisticum.org/ori.html}{Latin text}; accessed
2026-07-20.  Used for the tension between subsistence limits and inherited
anti-Jewish civil premises.
\end{itemize}

\subsection{Official Catholic reception}

\begin{itemize}
\item John XXII.  \emph{Redemptionem misit}, canonization of Thomas Aquinas,
18 July 1323, text and surrounding narratives in \emph{Fontes vitae}, fascicle
5.  Used only for the canonization act and the qualified traditions the bull
actually reports.

\item Pius V.  \emph{Mirabilis Deus}, 11 April 1567.  Doctor-of-the-Church
date and status controlled also by Paul VI's exact retrospective record below;
no broader doctrinal effect is inferred.

\item Leo XIII.  \emph{Aeterni Patris}, 4 August 1879,
\sourceurl{https://www.vatican.va/content/leo-xiii/en/encyclicals/documents/hf_l-xiii_enc_04081879_aeterni-patris.html}{Holy See English}; and
\emph{Cum hoc sit}, 4 August 1880,
\sourceurl{https://www.vatican.va/content/leo-xiii/it/briefs/documents/hf_l-xiii_brief_18800804_cum-hoc-sit.html}{Holy See Italian}; accessed
2026-07-20.  Used for the Thomistic renewal and patronage acts, not as a
critical biography.

\item Paul VI.  \emph{Lumen Ecclesiae}, 20 November 1974,
\sourceurl{https://www.vatican.va/content/paul-vi/it/letters/1974/documents/hf_p-vi_let_19741120_lumen-ecclesiae.html}{Holy See Italian}; accessed
2026-07-20.  Used for canonization and Doctor dates, the reception of local
censures, Thomistic authority, and the stated legitimacy of other schools.

\item Benedict XVI.  General Audiences on Thomas Aquinas, 2, 16, and 23 June
2010:
\sourceurl{https://www.vatican.va/content/benedict-xvi/en/audiences/2010/documents/hf_ben-xvi_aud_20100602.html}{life and work},
\sourceurl{https://www.vatican.va/content/benedict-xvi/en/audiences/2010/documents/hf_ben-xvi_aud_20100616.html}{faith and reason}, and
\sourceurl{https://www.vatican.va/content/benedict-xvi/en/audiences/2010/documents/hf_ben-xvi_aud_20100623.html}{moral and sacramental teaching}; accessed
2026-07-20.  Official catechesis is paraphrased; its retold anecdotes remain
tradition, not contemporary evidence.

\item Leo XIV.  Homily for the Jubilee of the World of Education, 1 November
2025,
\sourceurl{https://www.vatican.va/content/leo-xiv/en/homilies/2025/documents/20251101-messa-giubileo-formatori.html}{Holy See English}; accessed
2026-07-20.  Used for the current co-patronage formulation with John Henry
Newman.

\item Second Vatican Council.  \emph{Dignitatis humanae} 2--3, 10--12,
\sourceurl{https://www.vatican.va/archive/hist_councils/ii_vatican_council/documents/vat-ii_decl_19651207_dignitatis-humanae_en.html}{Holy See English};
\emph{Nostra aetate} 4,
\sourceurl{https://www.vatican.va/archive/hist_councils/ii_vatican_council/documents/vat-ii_decl_19651028_nostra-aetate_en.html}{Holy See English}; and
\emph{Gaudium et spes} 27, 29,
\sourceurl{https://www.vatican.va/archive/hist_councils/ii_vatican_council/documents/vat-ii_const_19651207_gaudium-et-spes_en.html}{Holy See English}; accessed
2026-07-20.  Used as later Catholic reception, never as wording attributed to
Thomas.

\item Catechism of the Catholic Church 369, 372, 2334--2335, and 2414,
\sourceurl{https://www.vatican.va/archive/ENG0015/_INDEX.HTM}{Holy See English
edition}; accessed 2026-07-20.  Used only for current Catholic reception on
equal dignity and slavery.  The following 2018 record, rather than this legacy
index, controls the revised wording of no.~2267.

\item Congregation for the Doctrine of the Faith.  ``New Revision of Number
2267 of the Catechism of the Catholic Church on the Death Penalty---Rescriptum
\emph{ex Audientia SS.mi},'' 1 August 2018,
\sourceurl{https://www.vatican.va/roman_curia/congregations/cfaith/documents/rc_con_cfaith_doc_20180801_catechismo-penadimorte_en.html}{Holy See English};
accessed 2026-07-20.  Used for the current wording of Catechism 2267, not as a
claim that Thomas anticipated it.

\item Catholic Church.  \emph{Calendarium Romanum}, editio typica (Vatican
City: Typis Polyglottis Vaticanis, 1969), p.~86; checked against the
\sourceurl{https://www.cultodivino.va/en/formazione/pubblicazioni/libri-liturgici/aliae/calendarium-romanum.html}{Dicastery publication record and later variations} and
\sourceurl{https://press.vatican.va/content/salastampa/en/bollettino/pubblico/2026/02/09/260209d.html}{current Holy See usage}; accessed 2026-07-20.  Used only for the 28 January
memorial in the General Roman Calendar.  Dominican, local, historical, and
other-rite calendars require separate edition and jurisdiction checks.
\end{itemize}

\subsection{Modern critical works}

\begin{itemize}
\item Torrell, Jean-Pierre.  \emph{Saint Thomas Aquinas, Volume 1: The Person
and His Work}, 3rd ed., trans. Robert Royal (Washington, DC: Catholic
University of America Press, 2022),
\sourceurl{https://www.cuapress.org/9780813235608/saint-thomas-aquinas/}{publisher
record}.  Principal critical control for life and work chronology; protected
prose is paraphrased.

\item Porro, Pasquale.  \emph{Thomas Aquinas: A Historical and Philosophical
Profile}, trans. Joseph G. Trabbic and Roger W. Nutt (Washington, DC: Catholic
University of America Press, 2016),
\sourceurl{https://www.cuapress.org/9780813230108/thomas-aquinas/}{publisher
record}.  Used for institutional setting, work development, and philosophical
profile by paraphrase.

\item Pasnau, Robert.  ``Thomas Aquinas,'' \emph{Stanford Encyclopedia of
Philosophy}, substantive revision 2022,
\sourceurl{https://plato.stanford.edu/entries/aquinas/}{online article};
accessed 2026-07-20.  Used for cautious chronology, works, 1277 reception, and
the warning that abundant later anecdote exceeds Thomas's self-disclosure.

\item Legge, Dominic.  ``Thomas Aquinas: A Life Pursuing Wisdom,'' in
\emph{The New Cambridge Companion to Aquinas} (Cambridge University Press,
2022),
\sourceurl{https://www.cambridge.org/core/books/abs/new-cambridge-companion-to-aquinas/thomas-aquinas/77405DF9F5717C511D856866F1E144F7}{publisher record}.
Used for recent biographical chronology, including the wider birth-year debate.

\item Mews, Constant J., and Marika Räsänen.  ``The Translation of the Holy
Body of Thomas Aquinas,'' publication record and repository links through the
\sourceurl{https://research.utu.fi/converis/portal/detail/Publication/176394361?lang=en_GB}{University of Turku}; accessed 2026-07-20.  Used to resist a
simple relic-custody narrative.  See also Räsänen, \emph{Thomas Aquinas's
Relics as Focus for Conflict and Cult in the Late Middle Ages} (Amsterdam,
2017).

\item Couvent des Jacobins, Toulouse,
\sourceurl{https://jacobins.toulouse.fr/fr/}{official site}; accessed
2026-07-20.  Used only for present institutional custody and display, not as
independent forensic authentication.

\item Hood, John Y. B.  \emph{Aquinas and the Jews} (Philadelphia: University
of Pennsylvania Press, 1995),
\sourceurl{https://www.pennpress.org/9780812215236/aquinas-and-the-jews/}{publisher
record}; and Rowan Dorin, study of the Duchess of Brabant letter,
\sourceurl{https://www.persee.fr/doc/jds_0021-8103_2016_num_2_1_6381}{article
record}.  Used by paraphrase to control anti-Jewish law and manuscript context.

\item Cavedon, Matthew P.  ``Early Stirrings of Modern Liberty in the Thought
of St. Thomas Aquinas,''
\sourceurl{https://www.cambridge.org/core/journals/politics-and-religion/article/early-stirrings-of-modern-liberty-in-the-thought-of-st-thomas-aquinas/767553BEC40952984E8DB03F49EA1DBC}{open-access article}; and Arthur Stephen
McGrade, ``The Medieval Idea of Heresy,''
\sourceurl{https://www.cambridge.org/core/journals/studies-in-church-history-subsidia/article/abs/medieval-idea-of-heresy-what-are-we-to-make-of-it/8F480C2841F772DFAD925010DA0DB452}{article record}.  Used to situate coercion
without denying Thomas's own conclusion.

\item Cornish, Paul J.  ``Marriage, Slavery, and Natural Rights in the
Political Thought of Aquinas,''
\sourceurl{https://www.cambridge.org/core/journals/review-of-politics/article/abs/marriage-slavery-and-natural-rights-in-the-political-thought-of-aquinas/FA41967FD9B128565CB404A66879187B}{article record}; and Boris O. Saavedra
Pérez, study of slavery and natural law,
\sourceurl{https://www.scielo.cl/scielo.php?pid=S0719-689X2019000100053}{open
article}; accessed 2026-07-20.  Used for the difference between absolute
natural slavery, \emph{ius gentium}, protected spheres, and abolitionism.

\item Newton, William.  ``Aquinas and the `Genius of Woman','' and Ana M.
Romero Carrasquillo and Mary L. Troyer de Romero, ``Aquinas on the Inferiority
of Woman,'' publisher records at
\sourceurl{https://journals.sagepub.com/doi/10.1177/0012580617730981}{SAGE} and
\sourceurl{https://www.pdcnet.org/acpq/content/acpq_2013_0087_0004_0685_0710}{Philosophy Documentation Center}; accessed 2026-07-20.  Used only to frame the
coexistence of rational equality, erroneous biology, and asserted hierarchy.
\end{itemize}
